\chapter{Metodologia}\label{chap:metodologia}

Este capítulo descreve o método que será empregado para alcançar os
objetivos propostos no \cref{chap:problema}. Inicialmente, a pesquisa é
caracterizada quanto ao seu propósito, abordagem e estilo. Em seguida,
são apresentados os materiais --- o escopo de culturas contemplado e as
tecnologias de implementação ---, a sequência de passos do método, a
estratégia de validação das hipóteses, o cronograma de execução, os
resultados esperados e as limitações do estudo.

%-------------------------------------------------------------------
\section{Caracterização da Pesquisa}\label{sec:caracterizacao}

Quanto ao propósito, esta pesquisa caracteriza-se como
\textbf{descritiva e aplicada}: descreve e estrutura um conjunto de
critérios técnicos já consolidados na literatura agronômica e aplica
esse conhecimento na construção e na avaliação de um artefato
computacional \citep{Wazlawick2014}. Quanto à abordagem, a pesquisa é
\textbf{quantitativa}, uma vez que a validação se dará pela medição
objetiva da concordância entre as saídas geradas pelo sistema e valores
de referência calculados de forma independente, evitando-se métricas
subjetivas.

Quanto ao estilo de pesquisa corrente em Computação, o trabalho
enquadra-se na \textbf{apresentação de um produto} (artefato de
software). O artefato a ser construído é um \textbf{sistema especialista
baseado em regras}, isto é, um sistema que reproduz o raciocínio de um
especialista humano em um domínio restrito por meio de uma base de
conhecimento explícita e de um motor de inferência que aplica regras do
tipo SE-ENTÃO \citep{Lanzer2010, SantosEtAl1994}. Conforme
\cite{Wazlawick2014}, o estilo de apresentação de produto exige cautela,
pois a simples demonstração de funcionamento possui baixa aceitação
científica: é necessário comparar o que foi construído com alguma
referência. Por essa razão, a validação do SIRAS não se limitará à
demonstração da aplicação: as saídas serão confrontadas com um oráculo
técnico --- os valores obtidos pela aplicação manual e independente dos
critérios do Manual de Calagem e Adubação \citep{ManualCalagem2016} ---
e o sistema será posicionado frente às ferramentas existentes
discutidas no \cref{chap:revisao}, das quais se diferencia por integrar
recomendação de manejo e estimativa de aptidão edáfica a partir de uma
análise individual de solo.

%-------------------------------------------------------------------
\section{Materiais}\label{sec:materiais}

Esta seção apresenta os materiais que delimitam e viabilizam o
desenvolvimento do SIRAS. Primeiro, define-se o \textbf{escopo de
culturas} contempladas, derivado diretamente do Manual de Calagem e
Adubação para os estados do RS e SC, que constitui a principal fonte de
conhecimento do sistema. Em seguida, descreve-se a \textbf{tecnologia de
implementação} adotada --- linguagem, framework, formato da base de
conhecimento e forma de distribuição --- e as \textbf{alternativas
tecnológicas} consideradas. Esses elementos correspondem, na metodologia
de pesquisa em Computação, aos materiais e recursos necessários à
execução do trabalho \citep{Wazlawick2014}.

\subsection{Escopo de Culturas}\label{sec:escopo}

O sistema abrangerá as culturas para as quais o Manual de Calagem e
Adubação \citep{ManualCalagem2016} disponibiliza, em seu capítulo de
recomendações, tabelas completas de adubação, totalizando
\textbf{61 culturas e grupos de culturas} organizados em seis grupos,
conforme detalhado a seguir. A contagem segue a própria organização do
Manual: cada item corresponde a uma seção de recomendação, de modo que
culturas que compartilham a mesma tabela --- como melancia e melão ---
são contabilizadas como uma única entrada. Com exceção da erva-mate,
discutida adiante, todas as culturas contempladas possuem pH de
referência definido na Tabela~5.1 do Manual, requisito do módulo de
calagem do sistema. Ficam fora do escopo: (i)~as culturas sem pH de
referência cuja lógica de calagem é incompatível com o modelo
padronizado adotado, como o arroz irrigado (cujo pH se eleva
naturalmente pela inundação), a mandioca e as espécies florestais
tolerantes à acidez; e (ii)~por delimitação de escopo, os grupos do
Manual voltados a forrageiras, palmeiras (juçara, real e pupunheira),
plantas medicinais, aromáticas e condimentares, plantas ornamentais e
sistemas especiais de produção, por estarem fora do foco de lavouras
comerciais do público-alvo do sistema.

\textbf{Culturas de grãos (21):} arroz de sequeiro, amendoim, aveia branca,
aveia preta, canola, centeio, cevada, ervilha seca e ervilha forrageira,
ervilhaca, feijão, girassol, linho, milho, milho pipoca, nabo forrageiro,
painço, soja, sorgo, tremoço, trigo e triticale. Esse grupo segue lógica
uniforme de recomendação, com tabelas de N, P e K baseadas na expectativa
de rendimento e na classe de disponibilidade dos nutrientes no solo; nas
leguminosas como soja, ervilhaca e tremoço, a adubação nitrogenada não é
recomendada em razão da fixação biológica de nitrogênio
\citep{ManualCalagem2016}.

\textbf{Hortaliças (18):} abóbora, abobrinha e moranga, alcachofra, alface,
almeirão, chicória, rúcula e salsa, alho, aspargo, berinjela, beterraba e cenoura,
brócolis e couve-flor, cebola, chuchu, ervilha, mandioquinha-salsa,
melancia e melão, nabo e rabanete, pepino, pimentão, repolho e tomate.

\textbf{Tubérculos (2):} batata e batata-doce, com tabelas completas de
N, P e K.

\textbf{Frutíferas (17):} abacateiro, ameixeira, amoreira-preta, bananeira,
caquizeiro, citros, figueira, macieira, maracujazeiro, mirtileiro, morangueiro,
nogueira-pecã, oliveira, pessegueiro e nectarineira, pereira, quivizeiro e videira.
Esse grupo apresenta complexidade adicional em relação aos demais: a adubação
é estruturada em até três fases (pré-plantio, crescimento e manutenção/produção),
e o sistema deverá adaptar as recomendações de acordo com a fase informada pelo
usuário. O fósforo e o potássio de pré-plantio seguem uma tabela compartilhada
entre todas as frutíferas (Tabela 6.5.1 do Manual), enquanto as doses de
nitrogênio e as recomendações de manutenção são específicas por espécie.

\textbf{Erva-mate (1):} cultura de expressiva relevância socioeconômica no
norte do Rio Grande do Sul. Constitui a única exceção ao critério de pH de
referência: o Manual a classifica entre as espécies tolerantes à acidez,
sem pH de referência definido, indicando a calagem apenas para o adequado
suprimento de cálcio e magnésio \citep{ManualCalagem2016}; o sistema
adotará esse tratamento diferenciado, sem cálculo de necessidade de
calagem por pH alvo. Apresenta a maior complexidade de implementação
entre as culturas contempladas: a adubação pode ocorrer em programa desde
o plantio (três fases: plantio/crescimento, formação de copa e produção) ou
em programa de recuperação de ervais já estabelecidos. A adubação de produção
é calculada com base na massa verde colhida (t/ha) e no tipo de manejo
(galho grosso retido ou retirado da área), variáveis que deverão ser informadas
pelo usuário ao selecionar essa cultura \citep{ManualCalagem2016}.

\textbf{Outras culturas comerciais (2):} cana-de-açúcar e tabaco, conforme
o agrupamento do próprio Manual. O \textbf{tabaco} é cultura de grande
importância econômica no RS e SC; as recomendações diferenciam-se entre os
tipos Virgínia e Burley, especialmente nas doses de nitrogênio e potássio,
e o sistema solicitará ao usuário o tipo cultivado ao selecionar essa
cultura. A \textbf{cana-de-açúcar} possui recomendações distintas para
cana-planta e cana-soca, com adubação nitrogenada definida pelo teor de
matéria orgânica do solo. Para ambas as culturas, o Manual estabelece pH
de referência de 6,0, com a particularidade de que a calagem é indicada
apenas quando o pH do solo for inferior a 5,5, regra que será incorporada
às regras de decisão do sistema \citep{ManualCalagem2016}.

\subsection{Tecnologia de Implementação}\label{sec:tecnologia}

O sistema será desenvolvido como uma aplicação web leve, executável
localmente em navegador, utilizando \textbf{Python}\footnote{Disponível
em: \url{https://www.python.org}.} como linguagem principal
com o microframework \textbf{Flask}\footnote{Disponível em:
\url{https://flask.palletsprojects.com}.} para a camada web, complementado por
\textbf{HTML} e \textbf{CSS} para a interface com o usuário. Essa escolha
fundamenta-se em três critérios principais: (i)~legibilidade e manutenibilidade
do código, uma vez que a lógica de regras agronômicas SE-ENTÃO em Python
resulta em código mais claro e auditável do que em linguagens de tipagem fraca;
(ii)~flexibilidade de estilização, pois o uso de HTML e CSS na interface permite
customização visual completa do formulário de entrada e do laudo de saída,
sem as limitações impostas por bibliotecas de interface nativas; e
(iii)~viabilidade para execução local, já que o Flask permite servir a aplicação
no próprio computador do usuário sem necessidade de servidor remoto ou conexão
com a internet, bastando executar o comando \texttt{python app.py} e acessar
o endereço \texttt{localhost} no navegador.

Está previsto o uso do Python na versão 3.12 (ou superior) e do Flask na
série 3.x, ambos software livre distribuído sob licenças permissivas
(PSF e BSD, respectivamente), compatíveis com o uso acadêmico e executáveis
em sistemas operacionais Windows ou GNU/Linux. Não há dependência de
serviços pagos, de banco de dados dedicado ou de conexão com a internet,
o que assegura que todos os recursos planejados estão disponíveis para a
execução do trabalho.

A lógica central do sistema --- o motor de inferência baseado em regras
SE-ENTÃO --- será implementada diretamente em Python, sem o uso de bibliotecas
externas de inferência (como Experta ou PyKE), de modo a manter o código
auditável e o comportamento das regras totalmente rastreável. As tabelas técnicas
do Manual de Calagem e Adubação \citep{ManualCalagem2016} serão armazenadas em
arquivos no formato JSON, carregados localmente pela aplicação. Essa separação
entre base de conhecimento e mecanismo de inferência é coerente com o paradigma
de sistemas especialistas descrito por \cite{Lanzer2010} e \cite{SantosEtAl1994},
e permite que as tabelas de recomendação sejam atualizadas de forma independente
da lógica de processamento. Para a geração do laudo de saída, será utilizada
a capacidade de renderização do navegador, com possibilidade de impressão
direta ou exportação em PDF pelo recurso nativo dos navegadores modernos.

Quanto à forma de distribuição, o SIRAS será disponibilizado como
\textbf{software de código aberto}, com o repositório publicado em
plataforma pública de versionamento. A forma de acesso prevista nesta
proposta é a \textbf{execução local}: o usuário obtém o código-fonte,
instala as dependências (Python e Flask) e executa a aplicação em sua
própria máquina, acessando a interface pelo navegador no endereço
\texttt{localhost}. Essa opção foi escolhida por dispensar custos de
hospedagem e por manter o controle dos dados na máquina do usuário, o
que é adequado ao escopo de uma proposta de TCC. A arquitetura web
adotada, contudo, não impede uma futura implantação em servidor remoto
--- bastaria hospedar a mesma aplicação Flask para que fosse acessada
diretamente por URL, sem instalação ---, possibilidade registrada como
trabalho futuro na \cref{sec:limitacoes}.

A \cref{fig:arquitetura} apresenta uma visão conceitual da arquitetura do sistema,
ilustrando os três módulos principais e o fluxo de processamento das informações.

\begin{figure}[htb!]
\centering
\begin{tikzpicture}[
  node distance = 0.6cm and 1.1cm,
  box/.style = {
    rectangle, rounded corners=6pt,
    minimum width=3.6cm, minimum height=1.0cm,
    text centered, font=\small\sffamily,
    line width=0.8pt
  },
  subbox/.style = {
    rectangle, rounded corners=3pt,
    minimum width=3.4cm,
    text centered, font=\footnotesize\sffamily,
    inner sep=4pt, line width=0.5pt
  },
  arrow/.style = {-{Stealth[length=7pt]}, line width=1.2pt},
  groupbox/.style = {
    rectangle, rounded corners=8pt,
    draw=gray!60, fill=gray!6,
    inner sep=8pt, line width=0.6pt
  }
]
%--- MÓDULO 1: ENTRADA DE DADOS ---
\node[box, fill=blue!15, draw=blue!50] (entrada) {Entrada de Dados};
\node[subbox, fill=blue!7, draw=blue!30, below=0.25cm of entrada] (e1) {pH, CTC, V\%, Al$^{3+}$};
\node[subbox, fill=blue!7, draw=blue!30, below=0.1cm of e1]       (e2) {P, K, MO, Textura};
\node[subbox, fill=blue!7, draw=blue!30, below=0.1cm of e2]       (e3) {Cultura selecionada};
%--- MÓDULO 2: MOTOR DE INFERÊNCIA ---
\node[box, fill=green!15, draw=green!50,
      right=1.1cm of entrada] (motor) {Motor de Inferência};
\node[subbox, fill=green!7, draw=green!30, below=0.25cm of motor] (m1) {Regras SE-ENTÃO};
\node[subbox, fill=green!7, draw=green!30, below=0.1cm of m1]     (m2) {Base de Conhecimento};
\node[subbox, fill=green!7, draw=green!30, below=0.1cm of m2]     (m3) {Manual RS/SC 2016};
%--- MÓDULO 3: LAUDO DE SAÍDA ---
\node[box, fill=orange!15, draw=orange!50,
      right=1.1cm of motor] (saida) {Laudo de Sa\'ida};
\node[subbox, fill=orange!7, draw=orange!30, below=0.25cm of saida] (s1) {Recomenda\c{c}\~ao de};
\node[subbox, fill=orange!7, draw=orange!30, below=0.1cm of s1]     (s2) {calagem e aduba\c{c}\~ao};
\node[subbox, fill=orange!7, draw=orange!30, below=0.1cm of s2]     (s3) {Aptid\~ao agr\'icola};
%--- CAIXAS DE AGRUPAMENTO ---
\begin{scope}[on background layer]
  \node[groupbox, fit=(entrada)(e1)(e2)(e3)] (g1) {};
  \node[groupbox, fit=(motor)(m1)(m2)(m3)]   (g2) {};
  \node[groupbox, fit=(saida)(s1)(s2)(s3)]   (g3) {};
\end{scope}
%--- SETAS ---
\draw[arrow, blue!60]   (g1.east) -- (g2.west);
\draw[arrow, green!60]  (g2.east) -- (g3.west);
\end{tikzpicture}
\caption{\label{fig:arquitetura}Arquitetura conceitual do SIRAS.}
\end{figure}

Do ponto de vista do usuário, a interação com o sistema é direta e
organiza-se em cinco passos, ilustrados na \cref{fig:fluxo}. O usuário
--- técnico agrícola, agrônomo ou produtor --- seleciona a cultura de
interesse, informa em um formulário os resultados de sua análise de
solo, aciona o processamento e, em poucos instantes, recebe na própria
tela o laudo com a recomendação de calagem e adubação e a estimativa de
aptidão, podendo então imprimi-lo ou exportá-lo em PDF. Todo o fluxo
ocorre localmente no navegador, sem necessidade de cadastro ou de
conexão com a internet.

\begin{figure}[htb!]
\centering
\includegraphics[width=\textwidth]{fig/Fluxo_de_Uso_Usuário}
\caption{\label{fig:fluxo}Fluxo de uso do SIRAS na perspectiva do usuário.}
\end{figure}


\subsection{Alternativas Tecnológicas}\label{sec:alternativas}

Caso durante o desenvolvimento sejam identificadas limitações ou necessidades
não previstas nesta proposta, as seguintes alternativas tecnológicas poderão
ser adotadas em substituição ou complemento à stack principal:

\begin{itemize}

  \item \textbf{Django} --- framework web Python mais completo que o Flask,
    indicado caso o sistema evolua para necessitar de autenticação de usuários,
    banco de dados relacional ou administração de dados mais robusta. Mantém
    a linguagem Python, minimizando o reaprendizado;

  \item \textbf{HTML + CSS + JavaScript puro} --- caso seja identificada a
    necessidade de eliminar qualquer dependência de servidor, a lógica de
    recomendação pode ser portada integralmente para JavaScript, resultando
    em uma aplicação que executa diretamente no navegador sem necessidade
    de instalação de Python ou Flask;

  \item \textbf{Streamlit} --- biblioteca Python voltada para criação rápida
    de aplicações de dados com interface web, indicada caso a prioridade
    seja a agilidade no desenvolvimento da interface em detrimento da
    customização visual, pois gera interfaces funcionais com poucas linhas
    de código Python;

  \item \textbf{SQLite} --- banco de dados leve e baseado em arquivo, sem
    necessidade de servidor de banco de dados, indicado caso seja necessário
    persistir histórico de análises realizadas pelo usuário entre sessões
    da aplicação.

\end{itemize}

%-------------------------------------------------------------------
\section{Método}\label{sec:metodo}

O método consiste na sequência de passos descrita a seguir, definida
a partir dos objetivos específicos enunciados na \cref{sec:obj_espec}.

\begin{enumerate}[label=\alph*),leftmargin=3em]

  \item \textbf{Estruturação da base de conhecimento:} extração das
    tabelas de interpretação e recomendação do Manual de Calagem e
    Adubação \citep{ManualCalagem2016} e sua transcrição para arquivos
    JSON, organizados pelos seis grupos de culturas definidos na
    \cref{sec:escopo}, com conferência manual dos valores transcritos;

  \item \textbf{Definição formal das regras:} especificação das regras
    SE-ENTÃO de interpretação dos atributos do solo (classes de
    disponibilidade), do cálculo da necessidade de calagem e do
    dimensionamento das doses de N, P\textsubscript{2}O\textsubscript{5}
    e K\textsubscript{2}O, incluindo os casos especiais das frutíferas
    (três fases), da erva-mate (programas de plantio e de recuperação)
    e do tabaco (tipos Virgínia e Burley);

  \item \textbf{Definição dos critérios de aptidão edáfica:} construção,
    a partir das faixas de interpretação dos atributos químicos e físicos
    do solo estabelecidas no Manual de Calagem e Adubação
    \citep{ManualCalagem2016} e das exigências agronômicas descritas na
    literatura de avaliação de aptidão das terras
    \citep{EmbrapaAptidao2025, Hofig2015}, das faixas de referência
    que sustentarão a classificação de aptidão por cultura, restrita aos
    atributos químicos e físicos do solo;

  \item \textbf{Implementação do módulo de entrada de dados:} formulário
    web para informar os resultados da análise de solo e a cultura de
    interesse, com validação de domínio das variáveis (faixas válidas,
    campos obrigatórios e unidades);

  \item \textbf{Implementação do motor de inferência:} desenvolvimento,
    em Python puro, do mecanismo que aplica as regras da base de
    conhecimento aos dados de entrada;

  \item \textbf{Implementação dos módulos de recomendação e de aptidão:}
    geração das recomendações de calagem e adubação e da estimativa de
    aptidão edáfica, integradas ao laudo de saída renderizado no
    navegador;

  \item \textbf{Construção do conjunto de casos de teste:} elaboração
    dos casos de validação e dos respectivos valores de referência,
    conforme detalhado na \cref{sec:validacao};

  \item \textbf{Validação e análise dos resultados:} execução dos casos
    de teste no sistema, apuração das taxas de concordância, comparação
    com ferramenta existente e decisão sobre as hipóteses formuladas na
    \cref{sec:hipoteses};

  \item \textbf{Avaliação de usabilidade:} aplicação do roteiro de
    tarefas e do questionário SUS com usuários do público-alvo, conforme
    a \cref{sec:val_usabilidade}, e coleta do retorno qualitativo;

  \item \textbf{Redação final:} consolidação da monografia, com a
    discussão dos resultados frente aos trabalhos relacionados, e
    revisão ortográfica e gramatical do texto.

\end{enumerate}

Está prevista a construção de um \textbf{protótipo incremental}, com dois
marcos intermediários de avaliação junto ao orientador: o primeiro
contempla o fluxo completo (entrada, inferência e laudo) restrito ao
grupo de culturas de grãos, cuja lógica de recomendação é uniforme; o
segundo estende o protótipo aos demais grupos, incorporando os casos de
maior complexidade (frutíferas, erva-mate e tabaco). Essa estratégia
permite validar a arquitetura cedo e reduzir o risco de retrabalho.

%-------------------------------------------------------------------
\section{Validação}\label{sec:validacao}

A validação definirá se as hipóteses formuladas na \cref{sec:hipoteses}
serão aceitas ou refutadas. Reconhece-se que a verificação da
concordância com o Manual atesta, sobretudo, que o sistema foi
\emph{implementado corretamente}; por isso, a estratégia de validação é
organizada em três dimensões complementares: (i)~\textbf{correção}, que
verifica a conformidade das saídas com os critérios técnicos de
referência; (ii)~\textbf{comparação}, que posiciona o SIRAS frente a uma
ferramenta existente; e (iii)~\textbf{usabilidade}, que avalia a
experiência de uso com usuários reais.

\subsection{Correção das saídas}\label{sec:val_correcao}

Esta dimensão verifica se as recomendações geradas estão corretas em
relação aos critérios técnicos vigentes. Para evitar métricas
subjetivas, adotam-se as seguintes \textbf{definições operacionais}:

\begin{itemize}

  \item \textbf{Concordância de recomendação:} um caso de teste é
    considerado concordante quando todas as saídas geradas pelo SIRAS
    --- necessidade de calagem (t/ha) e doses de N,
    P\textsubscript{2}O\textsubscript{5} e
    K\textsubscript{2}O (kg/ha) --- coincidem com os valores de
    referência do caso, admitida apenas a tolerância de arredondamento
    decorrente da precisão das tabelas do Manual;

  \item \textbf{Concordância de aptidão:} um caso de teste é considerado
    concordante quando a classe de aptidão edáfica atribuída pelo SIRAS
    coincide com a classe de referência derivada das faixas definidas
    na etapa~c) do método.

\end{itemize}

O conjunto de validação será composto por, no mínimo, \textbf{40 casos
de teste}, estratificados entre os seis grupos de culturas de modo a
cobrir as diferentes classes de interpretação dos atributos (teores
muito baixos a muito altos, diferentes CTCs e texturas) e os casos
especiais de maior complexidade. Os valores de referência de cada caso
constituirão o \textbf{oráculo} da validação e serão obtidos pela
aplicação manual e independente dos critérios do Manual, conferida com
o orientador. Para mitigar o risco de circularidade --- já que o
sistema é construído a partir do mesmo Manual usado na referência ---,
serão priorizados, quando disponíveis, laudos reais de análise de solo
acompanhados de recomendação técnica elaborada por profissional
habilitado, utilizados de forma anonimizada como casos de teste
adicionais.

Como o SIRAS é um sistema determinístico, execuções repetidas de um
mesmo caso produzem sempre a mesma saída; a análise dos resultados,
portanto, não requer testes de comparação de médias, e consistirá na
apuração da \textbf{proporção de casos concordantes} para cada hipótese,
acompanhada do respectivo intervalo de confiança de 95\%. A hipótese
H\textsubscript{1.1} será aceita se a taxa de concordância de
recomendação for igual ou superior a 80\%; a hipótese
H\textsubscript{1.2} será aceita se a taxa de concordância de aptidão
for igual ou superior a 80\%. Os casos discordantes serão analisados
individualmente e discutidos na monografia, distinguindo-se erros de
implementação (corrigíveis) de divergências de interpretação dos
critérios do Manual.

\subsection{Comparação com ferramenta existente}\label{sec:val_comparacao}

Para situar o SIRAS frente ao estado da prática, e seguindo a orientação
de que um artefato de software deve ser comparado a alguma referência
\citep{Wazlawick2014}, um subconjunto dos casos de teste será também
submetido a uma das ferramentas discutidas no \cref{chap:revisao}
--- como o FertFácil \citep{Tiecher2016} ou o AdubaTec
\citep{AdubaTec2021} ---, no que houver sobreposição de funcionalidade,
sobretudo na recomendação de calagem e adubação. Essa comparação é de
natureza \textbf{qualitativa}: pretende evidenciar convergências e
divergências entre as recomendações, documentar as funcionalidades que o
SIRAS oferece e que as ferramentas existentes não contemplam --- em
especial a estimativa de aptidão edáfica integrada à recomendação de
manejo --- e, assim, sustentar a contribuição do trabalho declarada na
H\textsubscript{1.2}. Eventuais divergências serão analisadas à luz das
diferenças de critério ou de versão de Manual adotadas por cada
ferramenta.

\subsection{Usabilidade}\label{sec:val_usabilidade}

A terceira dimensão avalia a experiência de uso do sistema com usuários
reais, aspecto não capturado pelas duas anteriores. Após a conclusão do
protótipo funcional, será conduzida uma avaliação de usabilidade com um
pequeno grupo de participantes do público-alvo (estudantes de agronomia,
técnicos agrícolas ou agrônomos), nos quais cada participante executará
um roteiro de tarefas típicas --- selecionar uma cultura, informar os
dados de uma análise de solo e interpretar o laudo gerado. Ao final, os
participantes responderão ao questionário \textbf{System Usability Scale
(SUS)} \citep{Brooke1996}, instrumento validado e amplamente utilizado
para mensurar a usabilidade percebida, que produz um escore de 0 a 100.
Como referência de interpretação, adota-se o valor consagrado na
literatura da escala, segundo o qual uma pontuação SUS média igual ou
superior a 68 indica usabilidade satisfatória \citep{Brooke1996}.
Complementarmente, será coletado retorno \textbf{qualitativo} por meio
de perguntas abertas sobre dificuldades encontradas e sugestões de
melhoria. Os resultados dessa etapa não testam as hipóteses
H\textsubscript{1.1} e H\textsubscript{1.2}, de natureza técnica, mas
fornecem evidências sobre a adequação da interface ao público-alvo e
subsidiam ajustes finais e recomendações de trabalhos futuros.

%-------------------------------------------------------------------
\section{Cronograma}\label{sec:cronograma}

A \cref{tab:cronograma} apresenta a distribuição das etapas do método
ao longo do período de execução do trabalho, em quinzenas, de agosto ao
início de dezembro de 2026, considerando que a apresentação do trabalho
está prevista para o início de dezembro. As letras correspondem às
etapas detalhadas na \cref{sec:metodo}, e os marcos de protótipo (M1 e
M2) estão indicados nas linhas correspondentes da tabela: M1, ao final
de setembro, corresponde ao fluxo completo restrito ao grupo de grãos;
M2, em meados de outubro, à extensão do protótipo aos demais grupos. A
redação da monografia ocorre de forma contínua a partir da segunda
quinzena de setembro; a validação técnica e a avaliação de usabilidade
concentram-se em novembro, e a primeira quinzena de dezembro fica
reservada à revisão final, entrega e preparação da apresentação.

\begin{table}[htb!]
\centering
\caption{\label{tab:cronograma}Cronograma de execução do trabalho (quinzenas, 2026).}
\small
\begin{tabular}{l c c c c c c c c c}
\hline
 & \multicolumn{2}{c}{\textbf{Ago}} & \multicolumn{2}{c}{\textbf{Set}} & \multicolumn{2}{c}{\textbf{Out}} & \multicolumn{2}{c}{\textbf{Nov}} & \textbf{Dez} \\
\textbf{Etapa} & 1 & 2 & 1 & 2 & 1 & 2 & 1 & 2 & 1 \\
\hline
a) Base de conhecimento (JSON)            & X & X &   &   &   &   &   &   &   \\
b) Regras SE-ENTÃO                        & X & X & X &   &   &   &   &   &   \\
c) Critérios de aptidão edáfica           &   & X & X &   &   &   &   &   &   \\
d) Módulo de entrada de dados             &   & X & X &   &   &   &   &   &   \\
e) Motor de inferência                    &   &   & X & X &   &   &   &   &   \\
\textbf{-- Marco M1 (protótipo grãos)}    &   &   &   & X &   &   &   &   &   \\
f) Módulos de recomendação e aptidão      &   &   &   & X & X &   &   &   &   \\
\textbf{-- Marco M2 (demais grupos)}      &   &   &   &   & X &   &   &   &   \\
g) Casos de teste e oráculo               &   &   &   & X & X & X &   &   &   \\
h) Validação e análise                    &   &   &   &   &   & X & X &   &   \\
i) Avaliação de usabilidade (SUS)         &   &   &   &   &   &   & X & X &   \\
j) Redação da monografia                  &   &   &   & X & X & X & X & X &   \\
k) Revisão final, entrega e apresentação  &   &   &   &   &   &   &   & X & X \\
\hline
\end{tabular}
\end{table}

%-------------------------------------------------------------------
\section{Resultados Esperados}\label{sec:resultados_esperados}

Caso os objetivos sejam alcançados, espera-se dispor de uma ferramenta
funcional, gratuita e executável localmente que apoie técnicos e
produtores do RS na interpretação de análises de solo, na obtenção de
recomendações de calagem e adubação e na avaliação da aptidão edáfica
de culturas específicas, reduzindo erros de interpretação e
qualificando a tomada de decisão. Espera-se, ainda, que a avaliação de
usabilidade indique boa aceitação da interface pelo público-alvo,
confirmando a adequação da ferramenta ao uso pretendido. No contexto do
crédito rural,
espera-se que o SIRAS possa servir de instrumento auxiliar de
conferência das recomendações técnicas que acompanham projetos
agrícolas, conferindo objetividade e rastreabilidade a essa análise.
Em termos de continuidade, a separação entre base de conhecimento e
motor de inferência deverá permitir a extensão futura do sistema a
outras culturas e a manuais de outras regiões, bem como sua utilização
como base para novos trabalhos acadêmicos.

%-------------------------------------------------------------------
\section{Limitações do Estudo}\label{sec:limitacoes}

O trabalho possui as seguintes limitações, que delimitam a abrangência
dos resultados:

\begin{itemize}

  \item a estimativa de aptidão é \textbf{exclusivamente edáfica},
    restrita aos atributos químicos e físicos informados na análise de
    solo, não contemplando fatores de relevo, clima, deficiência
    hídrica ou suscetibilidade à erosão considerados pelos sistemas de
    avaliação regional \citep{EmbrapaAptidao2025, Hofig2015};

  \item o escopo restringe-se às culturas com pH de referência e
    tabelas completas no Manual de Calagem e Adubação para os estados
    do RS e SC \citep{ManualCalagem2016}; culturas como arroz irrigado
    e mandioca, bem como critérios de outras regiões do país, estão
    fora do escopo desta proposta;

  \item a validação será realizada por casos de teste confrontados com
    valores de referência, e não por experimentos agronômicos de campo;
    portanto, os resultados atestam a conformidade do sistema com os
    critérios do Manual, e não o desempenho agronômico das
    recomendações em lavoura;

  \item o SIRAS é uma ferramenta de apoio à decisão e não substitui a
    responsabilidade técnica de profissional habilitado na emissão de
    recomendações oficiais de correção e adubação.

\end{itemize}

Registra-se, ainda, como possibilidade de \textbf{trabalho futuro}, a
implantação do sistema em servidor remoto, de modo a permitir o acesso
diretamente por URL, sem a necessidade de instalação local pelo usuário,
bem como a ampliação do escopo a outras regiões e a avaliação de
usabilidade com amostras maiores de participantes.
